\usepackage{threeparttable}
\usepackage{makecell}

\usepackage{longtable}
% Исправление ошибки longtable 4.24 при разбиении многостраничной таблицы.
% В longtable 4.27 этот фрагмент уже изменен, поэтому отсутствие старого
% шаблона не должно прерывать сборку.
\makeatletter
\patchcmd{\LT@output}
	{\vss}
	{\vskip\z@\@plus1fil\@minus\normalbaselineskip}
	{}
	{}
\patchcmd{\LT@output}
	{\vss}
	{\vskip\z@\@plus1fil\@minus\normalbaselineskip}
	{}
	{}
\makeatother
%\usepackage[none]{hyphenat}

\usepackage{rotating}
\usepackage{graphicx} % для resizebox


\usepackage{tabularx}
%\usepackage{cmap} % Улучшенный поиск русских слов в полученном pdf-файле
%\usepackage[T2A]{fontenc} % Поддержка русских букв
%\usepackage[utf8]{inputenc} % Кодировка utf8
\usepackage{pdfpages}
\usepackage{float}

% \usepackage[T2A]{fontenc}
% \usepackage[utf8]{inputenc}
% \usepackage[russian]{babel}
% \usepackage{enumitem}

\setlist[enumerate,1]{label={\arabic*)}}
% \setlist[enumerate,2]{label={\arabic{enumi}.\arabic{enumii})}}
\setlist[itemize]{label=--}

\usepackage[inkscapeformat=pdf]{svg}
\svgpath{/inc/img}

% \makeatletter
% \renewcommand*{\l@chapter}[2]{%
%   \ifnum\c@tocdepth>\m@ne
%     \addpenalty{-\@highpenalty}%
%     \vskip 1.0em \@plus\p@
%     \@dottedtocline{0}{0pt}{2.3em}{\bfseries #1}{\bfseries #2}%
%   \fi
% }
% \makeatother

\makeatletter
\renewcommand*{\l@chapter}[2]{
\ifnum \c@tocdepth>\m@ne
    \addpenalty{-\@highpenalty}
    \vskip 1em \@plus.2em
    \@dottedtocline{0}{0pt}{1.5em}{\bfseries #1}{\bfseries #2}
\fi
}
\makeatother

% \newcommand{\makeappendix}{%
%   \chapter*{\hfill{\centering ПРИЛОЖЕНИЕ А}\hfill}%
%   \addcontentsline{toc}{chapter}{ПРИЛОЖЕНИЕ А}%
% }

% Настройка заголовков
%\usepackage{titlesec}
%\makeatletter
%\renewcommand\LARGE{\@setfontsize\LARGE{22pt}{20}}
%\renewcommand\Large{\@setfontsize\Large{20pt}{20}}
%\renewcommand\large{\@setfontsize\large{16pt}{20}}
%\makeatother
%\titleformat{\chapter}{\large\bfseries}{\thechapter}{0.5em}{\large\bfseries}
%\titleformat{name=\chapter,numberless}[block]{\hspace{\parindent}}{}{0pt}{\large\bfseries\centering}
%\titleformat{\section}{\large\bfseries}{\thesection}{0.5em}{\large\bfseries}
%\titleformat{\subsection}{\large\bfseries}{\thesubsection}{0.5em}{\large\bfseries}
%\titleformat{\subsubsection}{\large\bfseries}{\thesubsection}{0.5em}{\large\bfseries}
%\titlespacing{\chapter}{12.5mm}{-22pt}{10pt}
%\titlespacing{\section}{12.5mm}{10pt}{10pt}
%\titlespacing{\subsection}{12.5mm}{10pt}{10pt}
%\titlespacing{\subsubsection}{12.5mm}{10pt}{10pt}

\let\oldsum\sum
\renewcommand{\sum}{\oldsum\limits}

\let\oldprod\prod
\renewcommand{\prod}{\oldprod\limits}


% Команда создания рисунка
\newcommand{\includesvgimage}[5]
{
	\ifthenelse{\equal{#2}{f}}
	{
		\begin{figure}[#3]
			\center{\includesvg[width=#4,inkscapelatex=false,extractformat={pdf,eps}]{inc/img/#1}}
			\caption{#5}
			\label{img:#1}
		\end{figure}
	}
	{
		\ifthenelse{\equal{#2}{w}}
		{
			\begin{wrapfigure}{#3}{#4}
				\center{\includesvg[width=#4,inkscapelatex=false,extractformat={pdf,eps}]{inc/img/#1}}
				\caption{#5}
				\label{img:#1}
			\end{wrapfigure}
		}
		{
			\PackageError{bmstu}{unknown image type}{}
		}
	}
}

\newcommand{\includepdfimage}[5]
{
	\ifthenelse{\equal{#2}{f}}
	{
		\begin{figure}[#3]
			\center{\includegraphics[width=#4]{inc/img/#1}}
			\caption{#5}
			\label{img:#1}
		\end{figure}
	}
	{
		\ifthenelse{\equal{#2}{w}}
		{
			\begin{wrapfigure}{#3}{#4}
				\center{\includegraphics[width=#4]{inc/img/#1}}
				\caption{#5}
				\label{img:#1}
			\end{wrapfigure}
		}
		{
			\PackageError{bmstu}{unknown image type}{}
		}
	}
}
